
%%%%%%%%%%%%%%%%%%%%%%%%%%%%%%%%%%%
%%%%%%%%%%%%%%%%%%%%%%%%%%%%%%%%%%%
%%%%%%%%%%%%%%%%%%%%%%%%%%%%%%%%%%%
%%%%%%%%%%%%%%%%%%%%%%%%%%%%%%%%%%%

\newpage

\thispagestyle{empty}
\vspace*{-5em} 

\renewcommand{\arraystretch}{1.2}
\begin{center}
    \textbf{IIS di Esempio - Comune Esempio (XX)}\\
    \textbf{Griglia di Valutazione
    \\Elaborato Scritto di Matematica e di Fisica - Liceo Artistico}
\end{center}
\begin{tabularx}{\textwidth}{|m{6.8cm}|m{5.6cm}|m{3.5cm}|}
\hline
\textbf{DOCENTE:} Cognome Nome &
\multicolumn{2}{|l|}{\textbf{DATA DELLA VERIFICA:} \rule{1.15cm}{0.4pt}/\rule{1.15cm}{0.4pt}/\rule{1.25cm}{0.4pt}}\\
\hline
\textbf{COGNOME:} &
\textbf{NOME:} &
\textbf{CLASSE:} \\
\hline
\end{tabularx}

\vspace{1mm}

\begin{tabularx}{\textwidth}{|>{\centering}m{3.3cm}|>{\centering\arraybackslash}m{11.8cm}|>{\centering\arraybackslash}m{0.8cm}|}
\hline
\textbf{INDICATORI} & \textbf{DESCRITTORI} & \textbf{PT}\\
\hline
\renewcommand{\arraystretch}{0.5}
&&\\
\end{tabularx}

\vspace{-0.75cm} % Spazio negativo per collegare le due tabelle

\renewcommand{\arraystretch}{1}
\begin{tabularx}{\textwidth}{|c|>{\arraybackslash}m{11.8cm}|>{\left\arraybackslash}m{0.8cm}|}

\multirow{7}{*}{\parbox{3.3cm}{ \textbf{\vspace{0.5cm}\\CONOSCENZE} \vspace{0.5cm}\\ \fontsize{9}{2}\selectfont{Conoscenza di\\ principi teorie,\\concetti, termini,\\regole, procedure,\\metodi e tecniche.}}}
& \fontsize{9}{2}\selectfont\textbf{{Nessuna} conoscenza}  dell’argomento. & 0 \\
\cline{2-3}
& \fontsize{9}{2}\selectfont{\textbf{Importanti lacune} nelle conoscenze.} & 1 \\
\cline{2-3}
& \fontsize{9}{2}\selectfont{Conoscenze \textbf{frammentarie e/o superficiali}.} & 2 \\
\cline{2-3}
& \fontsize{8.5}{2}\selectfont{Conoscenza della \textbf{maggior parte} degli argomenti teorici sebbene \textbf{con qualche incertezza}.} & 2,5 \\
\cline{2-3}
& \fontsize{9}{2}\selectfont{Conoscenza degli aspetti teorici \textbf{principali}.} & 3 \\
\cline{2-3}
& \fontsize{8}{2}\selectfont{Conoscenza \textbf{completa} degli aspetti teorici con \textbf{lievi incertezze} su conoscenze \textbf{non fondamentali}.} & 4 \\
\cline{2-3}
& \fontsize{9}{2}\selectfont{Conoscenze \textbf{esaurienti su tutti} gli aspetti teorici.} & \textbf{4,5} \\
\Xhline{2pt}

\multirow{7}{*}{\parbox{3.3cm}{ \textbf{\fontsize{10}{2}\selectfont{\vspace{0.1cm}\\CAPACIT\'{A} LOGICHE}} \vspace{0.5cm}\\ \fontsize{9}{2}\selectfont{Organizzazione e \\utilizzazione di\\conoscenze e abilità \\nella scelta di procedure.}}}
& \fontsize{9}{2}\selectfont{\textbf{Non individua} le procedure per la risoluzione.} & 0 \\
\cline{2-3}
& \fontsize{9}{2}\selectfont{Individuate procedure \textbf{sempre inadeguate} per la risoluzione dei problemi.} & 1 \\
\cline{2-3}
& \fontsize{9}{2}\selectfont{Individua \textbf{raramente} strategie  adeguate della prova.} & 1,5 \\
\cline{2-3}
& \fontsize{9}{2}\selectfont{Individua \textbf{solo  in parte}  le principali strategie  adeguate  per la risoluzione della prova.} & 2 \\
\cline{2-3}
& \fontsize{9}{2}\selectfont{Individua  \textbf{la  maggior parte} delle strategie  adeguate  per la risoluzione della prova.} & 2,5 \\
\cline{2-3}
& \fontsize{9}{2}\selectfont{Individua  \textbf{quasi  tutte} le procedure per la risoluzione della prova.} & 3,5 \\
\cline{2-3}
& \fontsize{9}{2}\selectfont{Individua  le procedure per la risoluzione della prova in maniera completa.} & \textbf{4,5} \\
\Xhline{2pt}

\multirow{4}{*}{\parbox{3.3cm}{ \textbf{\fontsize{8}{1}\selectfont{CAPACIT\'{A}\\ARGOMENTATIVA,\\\mbox{CHIAREZZA, LEGGIBILITÁ}\\E ORDINE}}}}
& \fontsize{9}{2}\selectfont{nessuna.} & 0,5 \\
\cline{2-3}
& \fontsize{9}{2}\selectfont{frammentaria.} & 1 \\
\cline{2-3}
& \fontsize{9}{2}\selectfont{essenziale.} & 1,5 \\
\cline{2-3}
& \fontsize{9}{2}\selectfont{completa.} & \textbf{2} \\
\Xhline{2pt}

\multirow{7}{*}{\parbox{3.3cm}{ \textbf{\scriptsize{\vspace{0cm}\\CORRETTEZZA\\\\ \mbox{DEGLI SVOLGIMENTI}}}\vspace{0.5cm}\\ \fontsize{9}{2}\selectfont{Correttezza nei calcoli,\\ nell’applicazione di\\tecniche e procedure\\nelle rappresentazioni\\geometriche e grafiche.\vspace{0.2cm}\\\fontsize{7}{2}\selectfont{(Le procedure mancanti\\verranno considerate errate)}}}}
& \fontsize{9}{2}\selectfont{Svolgimento \textbf{arbitrario ed incoerente} oppure non valutabile.} & 0,5 \\
\cline{2-3}
& \fontsize{9}{2}\selectfont{Svolgimento con una \textbf{procedura  gravemente  errata o esigua}.} & 1 \\
\cline{2-3}
& \fontsize{9}{2}\selectfont{Svolgimento con \textbf{rilevanti errori di calcolo o procedura}.} & 1,5 \\
\cline{2-3}
& \fontsize{9}{2}\selectfont{Svolgimento con \textbf{rilevanti errori di calcolo} in una \textbf{procedura quasi corretta}.} & 2,5 \\
\cline{2-3}
& \fontsize{8.5}{2}\selectfont{Svolgimento con un numero non rilevante di \textbf{errori non gravi nei calcoli e nel procedimento}.} & 3,5 \\
\cline{2-3}
& \fontsize{8.5}{2}\selectfont{Svolgimento con numero esiguo di \textbf{errori non gravi di calcolo} in un \textbf{procedimento corretto}.} & 4,5 \\
\cline{2-3}
& \fontsize{9}{2}\selectfont{Svolgimento \textbf{privo di errori} nei calcoli e nel procedimento.} & \textbf{5} \\
\Xhline{2pt}

\multirow{5}{*}{\parbox{3.3cm}{ \textbf{\\COMPLETEZZA}}}
& \fontsize{9}{2}\selectfont{Non svolto o svolto in misura esigua.} & 1 \\
\cline{2-3}
& \fontsize{9}{2}\selectfont{Svolgimento sommario.} & 1,5 \\
\cline{2-3}
& \fontsize{9}{2}\selectfont{Svolgimento di una parte rilevante.} & 2,5 \\
\cline{2-3}
& \fontsize{9}{2}\selectfont{Svolto quasi completamente.} & 3,5 \\
\cline{2-3}
& \fontsize{9}{2}\selectfont{Svolto completamente.} & \textbf{4} \\
\Xhline{2pt}
\end{tabularx}

\vspace{0.2cm}
\hspace{14.2cm} TOTALE: \hspace{1cm}/20
\vspace{0.5cm}

Comune Esempio,\raisebox{-1mm}{\rule{1cm}{0.4pt}}/\raisebox{-1mm}{\rule{1cm}{0.4pt}}/2026
\hspace{0.5cm} Docente:\rule{5cm}{0.4pt}
\hspace{2cm} Voto: \hspace{1.8cm}/10